\documentclass[12pt]{article}
\usepackage[T1]{fontenc}
\usepackage[utf8]{inputenc}
\usepackage{lmodern}
\usepackage{textcomp}
\usepackage{enumitem}
\usepackage{geometry}
\usepackage{graphicx}
\usepackage{amsmath, amssymb}
\geometry{margin=1in}
\begin{document}
\begin{center}
Economics - K-12 Level Assessment 25 \
Total Marks: 40 \
Answer all questions.
\end{center}

\section*{Multiple Choice (10 marks)}
\begin{enumerate}[label=\arabic*.]
\item The term "menu costs" refers to
\begin{enumerate}[label=(\alph*)]
    \item less choices due to inflation.
    \item financial assets being worth less due to inflation.
    \item a la carte savings falling.
    \item resource misallocation due to inflation.
\end{enumerate}
\item Which of the following would lead to a decrease in the money supply?
\begin{enumerate}[label=(\alph*)]
    \item The FED lowers the discount rate.
    \item The FED sells government securities in the secondary market.
    \item The federal government spends less money.
    \item The FED lowers reserve requirements.
\end{enumerate}
\item A recessionary gap exists when the short-run equilibrium level of real GDP
\begin{enumerate}[label=(\alph*)]
    \item decreases over time
    \item equals the full-employment level of real GDP
    \item is above the full-employment level of real GDP
    \item is below the full-employment level of real GDP
\end{enumerate}
\item What will happen to the equilibrium price and the equilibrium quantity of good Z when the price of good X which is a close substitute for Z rises?
\begin{enumerate}[label=(\alph*)]
    \item The equilibrium price will rise and the equilibrium quantity will fall.
    \item The equilibrium price will fall and the equilibrium quantity will rise.
    \item The equilibrium price and the equilibrium quantity will both rise.
    \item The equilibrium price and the equilibrium quantity will both fall.
\end{enumerate}
\item Which of the following is true of money and financial markets?
\begin{enumerate}[label=(\alph*)]
    \item As the demand for bonds increases the interest rate increases.
    \item For a given money supply if nominal GDP increases the velocity of money decreases.
    \item When demand for stocks and bonds increases the asset demand for money falls.
    \item A macroeconomic recession increases the demand for loanable funds.
\end{enumerate}
\end{enumerate}

\section*{True / False (10 marks)}
\textit{Answer True or False}
\begin{enumerate}[label=\arabic*.]
\setcounter{enumi}{5}
\item True or False: Flash estimates of GDP do not require any revisions once they are published.
\item True or False: Some components of M$_{2}$, like savings accounts, are less liquid than the components of M$_{1}$, such as cash and demand deposits.
\item True or False: Keynesian economists believe that fiscal policy is more effective than monetary policy in stimulating the economy.
\item True or False: When the price of one good increases and the quantity demanded of another good increases, they are considered substitute goods.
\item True or False: A tax cut in a full employment economy will lower real output while raising the price level.
\end{enumerate}

\section*{Long Form Response (20 marks)}
\textit{Answer each question with a clear and complete explanation.}
\begin{enumerate}[label=\arabic*.]
\setcounter{enumi}{10}
\item Discuss the various methods that labor unions employ to increase wages for their members and analyze which strategies are most effective, explaining why certain approaches are not utilized.
\item Discuss the effects of a rightward shift in the aggregate demand curve on the equilibrium price level and quantity of output, assuming the aggregate supply curve is upward sloping.
\end{enumerate}

\vfill
\noindent\textit{End of Paper}
\end{document}